% Vista científica exacta materializada desde una rama condicional.
% Fuente: source/demostraciones_primarias/09_06h_atlas_familias_rev6.tex; rama: HMTIncludeAtlasStructural.
% No modifica el contenido matemático de la rama de origen.
\chapter[Nonadic classification of families]{Nonadic classification of trajectory families}
\label{chap:atlas-genealogia-masas}
\index{nonadic atlas}
\index{family of routes}
\index{multisection}
\index{null sector}

The atlas does not begin with a list of particles or a table of masses.
It begins with a finite oriented support and with the readers that preserve the
history of its routes.  Its function is to answer a question prior to
any metrological comparison: what classes of persistence can the
nonadic geometry support, and what information must be preserved to distinguish them?

The answer has three cardinals, but not three rival inventories:
\[
 81\ \text{cells},\qquad
 56\ \text{route families},\qquad
 13\ \text{family--level multisections}.
\]
The first number counts oriented positions; the second identifies the
positions that publish the same finite route signature; the third gathers the
positions that share a family and depth relative to the center. None
of them counts known masses, names of species, or rows of an
external comparison.

\section{The complete geometry of the \(81\) cells}

Let
\[
 I_9=\{1,\ldots,9\},
 \qquad
 \mathcal C_{81}=I_9\times I_9.
\]
The coordinates \((r,c)\) preserve the two directions of APP\@. The parity
of each coordinate defines four internal sectors:
\[
 \operatorname{fam}(r,c)=
 \begin{cases}
  \ell^- ,& r\equiv1,\ c\equiv1\pmod2,\\
  \nu ,& r\equiv1,\ c\equiv0\pmod2,\\
  d ,& r\equiv0,\ c\equiv1\pmod2,\\
  u ,& r\equiv0,\ c\equiv0\pmod2.
 \end{cases}
 \tag{A1}\label{eq:familias-paridad-atlas}
\]
Depth in the atlas is not added as an independent label.  It is the
Chebyshev distance to the center:
\[
 G(r,c)=\max\{|r-5|,|c-5|\}.
 \tag{A2}\label{eq:generacion-chebyshev-atlas}
\]
Therefore, family and level follow from the position itself.  In particular,
the center is not obtained by selecting a species from outside: it is the unique
cell of depth zero.

\begin{theorem}[Census of families and generations of the atlas]
\label{thm:censo-atlas-81}
The maps \eqref{eq:familias-paridad-atlas} and
\eqref{eq:generacion-chebyshev-atlas} decompose the atlas in a
disjoint manner as
\[
 81=45_{\rm leptones}+36_{\rm quarks}
   =25_{\ell^-}+20_{\nu}+20_{d}+16_{u}.
 \tag{A3}\label{eq:censo-81-cuatro-familias}
\]
Refinement by \(G\) yields exactly the thirteen multisections
\[
\begin{array}{c|c|c}
\text{family}&\text{levels}&\text{cardinals}\\ \hline
\ell^-&G0,G2,G4&1,8,16\\
\nu&G1,G2,G3,G4&2,4,6,8\\
d&G1,G2,G3,G4&2,4,6,8\\
u&G1,G3&4,12.
\end{array}
\tag{A4}\label{eq:trece-multisecciones-atlas}
\]
The cell \((5,5)\) is the sole component of
\(\mathfrak S_{\ell^-,0}\).
\end{theorem}

\begin{proof}
In \(I_9\) there are five odd coordinates and four even ones.  The four parity
classes therefore have cardinalities
\[
 5\cdot5=25,\qquad
 5\cdot4=20,\qquad
 4\cdot5=20,\qquad
 4\cdot4=16,
\]
which proves \eqref{eq:censo-81-cuatro-familias}; the first two
classes comprise \(45\) cells and the last two, \(36\).

For the refinement, we enlarge squares centered at \((5,5)\).
In the odd--odd class there is one cell at radius zero, \(3^2-1=8\)
in the shell of radius two, and \(5^2-3^2=16\) in the shell of radius four.
In the odd--even class, the rectangles accumulated up to radii
\(1,2,3,4\) contain, respectively,
\[
 1\cdot2=2,\qquad 3\cdot2=6,\qquad
 3\cdot4=12,\qquad5\cdot4=20
\]
cells; their successive differences are \(2,4,6,8\). The
even--odd class gives the same calculation with the coordinates exchanged. In the
even--even class, there are \(2^2=4\) cells at radius one and
\(4^2-2^2=12\) at radius three. These shells are disjoint and exhaust
\(I_9^2\), so no other multisection remains.
\end{proof}

The word \emph{level} denotes in this chapter the combinatorial depth
\(G\).  The subsequent physical reading preserves that geometry but does not
define it.  This prevents particle names from being turned into premises
of the census that must explain them.

\section{Central inversion, \(41\) orbits, and uniqueness of the centre}

The inverse of the atlas is
\[
 \rho_{\rm c}(r,c)=(10-r,10-c).
 \tag{A5}\label{eq:inversion-central-atlas-rev6}
\]
Since \(10-r\) has the same parity as \(r\), inversion preserves the
four families.  Moreover,
\[
 |(10-r)-5|=|r-5|,
 \qquad
 |(10-c)-5|=|c-5|,
\]
and therefore preserves the level \(G\).  Each of the thirteen
multisections is \(\rho_{\rm c}\)-stable.

\begin{proposition}[Orbits and Central Exception]
\label{prop:orbitas-centro-atlas-rev6}
The action of \(\rho_{\rm c}\) on \(\mathcal C_{81}\) has exactly
\(41\) orbits.  Its unique fixed point is \((5,5)\).  Consequently, the
multisection \(\mathfrak S_{\ell^-,0}\) admits the unique pointwise
\(\rho_{\rm c}\)-equivariant selector; in the other twelve, the natural object is the
orbit or the complete multisection.
\end{proposition}

\begin{proof}
The equation
\[
 (r,c)=\rho_{\rm c}(r,c)
\]
is equivalent to \(2r=2c=10\), hence \(r=c=5\).  The other eighty cells are
grouped into forty pairs, whence \(1+40=41\) orbits.

Let \(y\) be a noncentral multisection and suppose that an equivariant point selector
chooses \(a(y)\in\mathfrak S_y\).  Since \(\rho_{\rm c}\) acts
trivially on the family--level type, equivariance would impose
\[
 a(y)=a(\rho_{\rm c}y)=\rho_{\rm c}a(y).
\]
The chosen point would have to be fixed, but the only fixed point lies in
\(\mathfrak S_{\ell^-,0}\).  The contradiction proves the assertion.
\end{proof}

This proposition explains by which a species internal is not, in general, a
cell isolated.  The electron central is exceptional because its multisection
is reduceds to a point.  In the remaining families, the orientation conjugate form
part of the identity that must is preserved.

\section{The \(56\) route families}

Each cell bears the finite shadow of its transport history:
\[
 \mathcal R(r,c)=
 \bigl(
 n_A,\ s_C,\ k_{\Delta_4},\
 \nu_{120},\ \nu_{270},\ \tau_{12}
 \bigr)_{r,c}.
 \tag{A6}\label{eq:lector-familia-ruta-atlas}
\]
The first five coordinates are integers; \(\tau_{12}\) is the oriented
dodecaphase word. Two cells belong to the same route family
if and only if they have the same sextuple \(\mathcal R\).

\begin{theorem}[Census of route families and minimal memory]
\label{thm:censo-56-familias-ruta-rev6}
The image of \(\mathcal R\) contains exactly \(56\) elements.  The
histogram of the sizes of its fibres is
\[
 41\times1,\qquad
 10\times2,\qquad
 2\times3,\qquad
 1\times4,\qquad
 2\times5.
 \tag{A7}\label{eq:histograma-fibras-ruta}
\]
In particular,
\[
 41+10+2+1+2=56,
\]
\[
 41\cdot1+10\cdot2+2\cdot3+1\cdot4+2\cdot5=81.
 \tag{A8}\label{eq:doble-censo-rutas}
\]
Fixing the lexicographic order within each fiber, the rank
\[
 \kappa(r,c)\in\{0,1,2,3,4\}
\]
makes the application injective
\[
 (r,c)\longmapsto\bigl(\mathcal R(r,c),\kappa(r,c)\bigr).
 \tag{A9}\label{eq:lift-memoria-cinco-atlas}
\]
Five is the minimum size of a uniform alphabet capable of performing this lift.
\end{theorem}

\begin{proof}
The exhaustive enumeration of \(I_9^2\) groups the eighty-one sextuples
\eqref{eq:lector-familia-ruta-atlas} by coordinatewise equality.
It produces the five fiber sizes of
\eqref{eq:histograma-fibras-ruta}; the two identities
\eqref{eq:doble-censo-rutas} simultaneously verify the number of classes
and the coverage of the eighty-one cells.

The lexicographic rank distinguishes the elements of each fiber, so that
\eqref{eq:lift-memoria-cinco-atlas} is injective.  Since two fibers
of cardinal five exist, any common alphabet with fewer than five symbols
would identify two cells of one of them by the pigeonhole principle.
\end{proof}

The rank \(\kappa\) is positional memory within a fiber; it is not
charge, mass, or color. Its function is to recover the cell that the visible
route signature had identified with others. The number \(56\) counts classes of the
sextuple \(\mathcal R\), whereas the number \(13\) counts classes of the pair
\((\operatorname{fam},G)\). They are different quotients of the same atlas.

\paragraph{Cellular and orbital quotients of the atlas: domains and projection maps.}
The raw CT--108 projection and the scientific classification read different
data from the same eighty-one cells. If
\(q_{\rm CT}\) preserves only the form of the dodecaphase tritic word,
while \(\mathcal R\) preserves the route sextuple and
\(q=(\operatorname{fam},G)\) preserves family and depth, the typed
table is
\[
\begin{array}{rcll}
 \mathcal C_{81}
 &\xrightarrow{\ q_{\rm CT}\ }&
 \mathcal W^{\rm raw}_{7}
 & \text{raw CT--108 forms},\\[2mm]
 \mathcal C_{81}
 &\xrightarrow{\ \mathcal R\ }&
 \Sigma^{\rm ruta}_{56}
 & \text{route families},\\[2mm]
 \mathcal C_{81}
 &\xrightarrow{\ q\ }&
 \Sigma_{13}
 & \text{family--level multisections}.
\end{array}
\tag{A9a}\label{eq:dos-cocientes-atlas-rev8}
\]
The seven raw fibres have multiplicities
\[
 54,2,2,2,7,7,7,
 \qquad54+3\cdot2+3\cdot7=81.
\]
Therefore, the archaeological census \(81\to7\) does not replace the atlas \(81/56/13\). Nor does the abbreviated notation \(81/56/13\) assert the existence of a canonical arrow \(56\to13\): some fibers of \(\mathcal R\) traverse more than one family--level pair. The maps \(\mathcal R\) and \(q\) are parallel quotients and preserve different invariants. In the quark sector, the active convention remains
\[
 \operatorname{col}(r,c)=r-c\pmod3;
\]
the historical notation \(c-r\) is its conjugate convention, not another color or
another quotient of the atlas.

\section{Distinction between cell, route family, and multisection}

Define
\[
 q:\mathcal C_{81}\longrightarrow\Sigma_{13},
 \qquad
 q(r,c)=\bigl(\operatorname{fam}(r,c),G(r,c)\bigr),
 \tag{A10}\label{eq:cociente-trece-multisecciones}
\]
and, for \(y\in\Sigma_{13}\),
\[
 \mathfrak S_y^{\rm atlas}
 =q^{-1}(y)
 =\{(r,c)\in\mathcal C_{81}:q(r,c)=y\}.
 \tag{A11}\label{eq:multiseccion-atlas-rev6}
\]
The four levels of reading are then separated:

\begin{enumerate}[label=(\roman*)]
 \item a \emph{cell} is an oriented position \((r,c)\), with all its
 local data;
 \item a \emph{route family} is a fibre of \(\mathcal R\): it preserves
 the same finite transport signature;
 \item a \emph{multisection} is a fiber of \(q\): it preserves the family,
 depth, and all its conjugate orientations;
 \item a \emph{physical representation} is the subsequent realization of an
enriched multisection in a state space; it is not another partition of the
board or the arbitrary choice of one of its cells.
\end{enumerate}

This distinction gives mathematical content to the word \emph{species}. The internal type is the multisection. A concrete particle is a persistent section realized within it. The physical representation preserves the action, orientation, and observables of the type; it does not work backward to redefine the atlas from an already known mass.

\section{Color residual as ternary orientation}

In the quark sector, formed by the \(36\) cells with even \(r\), we define
\[
 \operatorname{col}(r,c)=r-c\pmod3,
 \tag{A12}\label{eq:color-residual-atlas-rev6}
\]
with the dictionary
\[
 0\longmapsto R,\qquad1\longmapsto G,\qquad2\longmapsto B.
\]

\begin{proposition}[Residual color balance]
\label{prop:color-atlas}
The thirty-six quark cells are distributed exactly as
\[
 36=12_R+12_G+12_B.
 \tag{A13}\label{eq:balance-color-atlas-rev6}
\]
The central inversion fixes \(R\) and interchanges \(G\) and \(B\).
\end{proposition}

\begin{proof}
Upon fixing any of the four even values of \(r\), as
\(c=1,\ldots,9\) is traversed, each class modulo three appears three times. Each residue
therefore has \(4\cdot3=12\) representatives. Moreover,
\[
 \operatorname{col}\bigl(\rho_{\rm c}(r,c)\bigr)
 =(10-r)-(10-c)
 =-(r-c)\pmod3.
\]
The modular negation fixes the zero class and swaps the one and two classes.
\end{proof}

The tritic \(R,G,B\) is therefore an exact observable of the atlas and not a
decoration added to a table of quarks.  The color representation is
subsequently constructed on this ternary support; the census \(12+12+12\) is its
prior geometry.

\section{From the surviving extension to the signature}

The atlas classifies finite cuts.  The identity of a particle requires the
deep history from which that cut proceeds.  Let
\[
 \Omega_{\rm HMT}^{\rm surv}
 =\varprojlim_m\mathcal S_m
\]
the inverse system of surviving states.  The projection onto the atlas \(\pi_{81}\) and the route
shadow form the chain
\[
 \Omega_{\rm HMT}^{\rm surv}
 \xrightarrow{\ \operatorname{sh}_{108}\ }
 \Gamma_{108}^{\rm enr}
 \xrightarrow{\ \mathcal E\ }
 \mathcal L_{108}
 \xrightarrow{\ L_{\rm tick}\ }
 \mathbb Z^5.
 \tag{A14}\label{eq:cadena-extension-firma-atlas-rev6}
\]
Here \(\operatorname{sh}_{108}\) publishes a route of \(108\) steps without
forgetting sheet, orientation, boundary, or predecessors; \(\mathcal E\) forms the
dodecaphase record; and \(L_{\rm tick}\) extracts
\[
 v(x)=
 \bigl(n_A,s_C,k_{\Delta_4},\nu_{120},\nu_{270}\bigr).
 \tag{A15}\label{eq:firma-cinco-extension-rev6}
\]
None of these three arrows refers to a mass, a unit, or the name of
a particle.

The reason for preserving the ledger before projection is apparent in its twelve
oriented events \(e_0,\ldots,e_{11}\). For \(0\leq i\leq5\), let
\[
 p_i=e_i+e_{i+6},
 \qquad
 a_i=e_i-e_{i+6},
 \qquad
 s_C=\sum_{i=0}^{5}p_i,
\]
and, for \(0\leq i<5\),
\[
 M_i=p_i-p_5.
\]
Then
\[
 p_5=\frac{s_C-\sum_{i=0}^{4}M_i}{6},
 \qquad
 p_i=p_5+M_i,
\]
\[
 e_i=\frac{p_i+a_i}{2},
 \qquad
 e_{i+6}=\frac{p_i-a_i}{2}.
 \tag{A16}\label{eq:reconstruccion-registro-doce}
\]
The identity
\[
 1+5+6=12
\]
is thus a reversible reconstruction: one total, five differences, and six
orientations recover the twelve events.  Projecting onto the five-integer
signature before composing histories would eliminate precisely the placement and
orientation that make it possible to decide whether two routes can be composed.

The quotient \(q\) is lifted to the deep history through projection to the atlas
\(\pi_{81}:\Omega_{\rm HMT}^{\rm surv}\to\mathcal C_{81}\):
\[
 \widetilde{\mathfrak S}_y
 =\{x\in\Omega_{\rm HMT}^{\rm surv}:q(\pi_{81}x)=y\}.
 \tag{A17}\label{eq:multiseccion-enriquecida-rev6}
\]
This is the enriched multisection.  Its shadow belongs to the atlas; its route
belongs to \(\Gamma_{108}^{\rm enr}\); its complete memory belongs to the
record; and its evaluation is applied only at the end.  The genealogy is not a
note appended to the value: it is the object that makes the value possible and preserves the
other readings of the same section.

\section{Null sector bifurcation prior to the positive character}

The mass character acts on the entire signature and has positive codomain:
\[
 \chi_{\rm mass}:\mathbb Z^5\longrightarrow\mathbb R_{>0}.
 \tag{A18}\label{eq:caracter-positivo-atlas-rev6}
\]
By definition, there is no signature \(v\) with
\(\chi_{\rm mass}(v)=0\). The null sector is, consequently, not a
fourteenth multisection nor the degenerate limit of the thirteen preceding ones.

Its algebraic support belongs to the split branch.
\[
 \mathcal A_{\rm sp}
 =\mathbb R[\mathsf R]/(\mathsf R^2-1),
 \qquad
 P_\pm=\frac{1\pm\mathsf R}{2}.
\]
Every element is written in a unique way as
\[
 z=r_+P_++r_-P_-,
 \qquad
 N_{\rm sp}(z)=r_+r_-,
\]
and the two lines
\[
 \mathbb R P_+,\qquad\mathbb R P_-
\]
are null because \(N_{\rm sp}(\lambda P_\pm)=0\). A null section is an
open route transported along one of those directions. The material interior
requires the coupling of both, for which \(r_+r_->0\).

Thus are typed two branches:
\[
 \text{material persistence}
 \longrightarrow\mathbb Z^5
 \xrightarrow{\chi_{\rm mass}}\mathbb R_{>0},
\]
\[
 \text{null propagation}
 \longrightarrow\partial_{\rm null}\mathcal C_{\rm sp}^{+}.
\]
The photonic and gauge realizations subsequently read the second branch.  Their
separation from the positive character prevents the photon and gluon from being erased because they do not appear
in a list of masses, while also preventing a null mass from being fabricated by means of a
character whose codomain excludes it.

\section{Finite classification system and path memory}

The three quotients of the chapter fulfill complementary functions:
\[
 \mathcal C_{81}
 \xrightarrow{\ \mathcal R\ }\mathcal R_{56},
 \qquad
 \mathcal C_{81}
 \xrightarrow{\ q\ }\Sigma_{13},
 \qquad
 \mathcal C_{81}
 \xrightarrow{\ /\langle\rho_{\rm c}\rangle\ }
 \mathcal O_{41}.
\]
The first preserves the route signature; the second preserves family and level;
the third preserves central conjugation.  Together with residual color and the
lift \(\kappa\), they form a finite system capable of reconstructing
which information is seen, which is identified, and which must remain as
memory.

The causal chain of the atlas is thereby fixed:
\[
 \boxed{
 \begin{gathered}
 \text{surviving extension}
 \longrightarrow\text{enriched route}
 \longrightarrow\text{record}
 \longrightarrow\text{signature}\\
 \longrightarrow\text{character}
 \longrightarrow\text{dimensional reading}
 \end{gathered}}
 \tag{A19}\label{eq:cadena-rectora-atlas-rev6}
\]
The eighty-one cells constitute the local domain; the fifty-six
families express route coincidences; the thirteen multisections are the
internal types; the forty-one orbits preserve conjugation; and the
center \((5,5)\) is the sole pointwise exception. Mass appears thereafter
as a reading of an already constructed persistence. The atlas is not an
abridged table of results: it is the structure explaining why a table can
exist.
